\documentclass[12pt,a4paper]{article}

\usepackage{iftex}
\usepackage{amsmath,amssymb}
\ifPDFTeX
  \usepackage[T2A]{fontenc}
  \usepackage[utf8]{inputenc}
\else
  % XeLaTeX / LuaLaTeX (в т.ч. tectonic): Unicode-шрифт с кириллицей
  \usepackage{fontspec}
  \setmainfont{cmunrm.otf}[BoldFont=cmunbx.otf, ItalicFont=cmunti.otf,
    BoldItalicFont=cmunbi.otf]
\fi
\usepackage[russian,english]{babel}
\usepackage[margin=2.5cm]{geometry}
\usepackage[hidelinks]{hyperref}
\usepackage{xcolor}

% Меньше вылезающих за поле строк с длинными латинскими терминами
\setlength{\emergencystretch}{2em}

% Пометки для черновика

\title{Детерминированные алгоритмы аппроксимации первых слоёв LLM}
\author{Чанкаев Шамиль, Мархабаев Бекет}
\date{}

\begin{document}
\selectlanguage{russian}
\maketitle

% =====================================================================
\begin{abstract}
% Аннотация (подключается в annotation.tex и main.tex)
Первые слои больших языковых моделей (LLM) на основе архитектуры
Transformer~\cite{vaswani2017attention} выполняют сравнительно простые
операции: восстанавливают статистику биграмм и коротких
$n$-грамм~\cite{elhage2021framework}, реализуют фиксированные позиционные
паттерны внимания и хранят поверхностные шаблоны в слоях
прямого распространения~\cite{geva2021ffn}.
Мы исследуем гипотезу о том, что вычисления первых $k$ слоёв можно
заменить \emph{детерминированными} алгоритмами, то есть процедурами без
случайной выборки и без дообучения градиентным спуском, которые
строятся по весам модели и статистике калибровочного корпуса в
замкнутой форме. Мы адаптируем к каузальному декодеру Qwen3 методы
эффективного внимания и сжатия слоёв MLP из литературы и предлагаем
собственные: факторизацию внимания <<позиционное ядро $\times$
значимость ключевого токена>>, типизированную замену голов,
линеаризацию MLP с компенсацией дрейфа остаточного потока и табличные
замены блоков. Эксперименты проводятся на Qwen3-0.6B и 1.7B
(Base) с оценкой по KL-дивергенции к исходной модели, perplexity и
стандартным бенчмаркам.
Профилирование показывает, что первые слои образуют отдельный режим
<<до стока внимания>>: массивная активация первого токена и сток
внимания возникают лишь после нескольких первых слоёв, что задаёт
естественную глубину замены $k_{\mathrm{sink}}$. Мы также показываем,
что рандомизированные признаки FAVOR+ (Performer~\cite{choromanski2021performer})
неприменимы к первым слоям без переобучения: нормы ключей после
QK-нормализации лежат далеко за пределами области, где их оценка
ядра устойчива.
Основные результаты. (1) Внимание первых трёх слоёв заменяется окном
из 64 позиций со стоком с ростом perplexity 1.5--1.8\,\% и без потерь на
семи бенчмарках; MLP первых слоёв заменяется хуже: лучший метод "---
предложенный детерминированный предиктор активных нейронов "--- стоит
4.7\,\% (0.6B) и 1.7\,\% (1.7B) при четверти вычислений MLP, а по карте
заменяемости MLP ранних слоёв оказываются самыми незаменимыми в модели.
(2) Мы обнаруживаем ловушку калибровки в замкнутой форме: единичные
токены с массивной активацией доминируют в задаче наименьших квадратов,
ошибка ложится в каналы стока внимания и обрушивает модель; мы даём
формулу влияния строки, воспроизводим эффект дозой отравленных окон и
устраняем его исключением строк-выбросов, а также показываем, что
массивную активацию создают один--пять нейронов, удаление которых
разрушает модель. (3) Априорные статистики слоя предсказывают KL его
замены с ранговой корреляцией 0.9--0.97 и позволяют планировать замены
на непросмотренной модели в 10--15 раз точнее замены префикса.
(4) На втором семействе (SmolLM2-1.7B, Llama-3.2-1B) сток оказывается
привязан не к позиции, а к токену: один нейрон создаёт массивную
активацию на <<the>> в любой позиции, и замены переносятся лишь с
защитой таких нейронов и хранением таких ключей (окно: $64\to4\,\%$ PPL
при $k=8$); MLP первых слоёв у этого семейства заменяется хуже, а
планировщик переносится.
Результаты дополняют дискуссию о том, насколько внутренние вычисления
обученных сетей <<уже чисты>> и поддаются явному
описанию~\cite{cammarata_clean,gao_notclean}.

\medskip
\noindent\textbf{Ключевые слова:} трансформеры, большие языковые модели,
аппроксимация внимания, механистическая интерпретируемость,
детерминированные алгоритмы, низкоранговые приближения, сжатие MLP.

\end{abstract}

% =====================================================================
% Введение (подключается в annotation.tex и main.tex)
\section{Введение}

Архитектура Transformer~\cite{vaswani2017attention} стала основой
современных больших языковых моделей. Её центральный механизм,
самовнимание (self-attention), требует $O(n^2)$ операций по длине
последовательности $n$. Это породило большое семейство методов
эффективной аппроксимации внимания. Их можно условно разделить на
несколько групп:
\begin{itemize}
  \item \textbf{ядерные и линейные} методы, в которых softmax заменяется
    скалярным произведением признаковых отображений: линейное
    внимание~\cite{katharopoulos2020rnn}, Performer со случайными
    признаками FAVOR+~\cite{choromanski2021performer}, случайные
    фурье-признаки RFA~\cite{peng2021rfa}, SOFT без
    softmax~\cite{lu2021soft};
  \item \textbf{низкоранговые} методы: проекции Linformer по оси
    последовательности~\cite{wang2020linformer}, приближение Нистрёма
    в Nyströmformer~\cite{xiong2021nystrom}, аналитический фреймворк
    A3~\cite{wong2025a3};
  \item \textbf{разреженные} шаблоны внимания: Sparse
    Transformer~\cite{child2019sparse}, хеширование LSH в
    Reformer~\cite{kitaev2020reformer}, комбинация локального,
    глобального и случайного внимания в BigBird~\cite{zaheer2020bigbird};
  \item \textbf{замена внимания} фиксированным или
    синтезированным смешиванием токенов: преобразование Фурье в
    FNet~\cite{leethorp2022fnet}, Synthesizer~\cite{tay2021synthesizer}.
\end{itemize}

Почти все эти методы аппроксимируют механизм внимания \emph{равномерно
по глубине}: одна и та же схема применяется во всех слоях и, как
правило, требует обучения модели с нуля или дообучения. При этом
многие из них по своей природе \emph{стохастичны}: случайные признаки,
случайные хеш-функции, случайные рёбра графа внимания. Из-за этого
ошибка аппроксимации зависит от выборки, а результат вычислений
воспроизводим лишь с точностью до зерна генератора.

Параллельно развивается механистическая интерпретируемость: она
показывает, что слои трансформера неоднородны по своей
функции~\cite{olah2020zoom,rai2024review}. Математический анализ
трансформерных схем~\cite{elhage2021framework} показывает, что модель
без слоёв внимания приближает статистику биграмм. Однослойная модель,
состоящая только из внимания, дополнительно реализует
<<skip-trigram>>-зависимости. Индукционные головы, которые
обеспечивают обучение в контексте, возникают лишь при композиции как
минимум двух слоёв внимания~\cite{olsson2022induction}. Слои прямого
распространения работают как память <<ключ--значение>>, причём ключи
нижних слоёв реагируют на поверхностные шаблоны (конкретные
$n$-граммы), а верхних "--- на семантические~\cite{geva2021ffn}.
Словарное обучение на активациях MLP небольшой модели выделяет
интерпретируемые моносемантические признаки~\cite{bricken2023monosemanticity}.
Известно также, что многоголовое самовнимание способно выражать
свёрточные слои~\cite{cordonnier2020relationship}. Значит, локальные
позиционные паттерны внимания в принципе описываются фиксированными
ядрами. Наконец, теоретический анализ
показывает, что чистое внимание без остаточных связей и MLP
вырождается по рангу дважды экспоненциально с
глубиной~\cite{dong2021attention}. Поэтому остаточный поток играет
ключевую роль при любой послойной замене вычислений.
Отдельно идёт дискуссия о том, являются ли обученные нейросети
<<уже чистыми>> (\emph{already clean}), то есть поддающимися явному
описанию~\cite{cammarata_clean,gao_notclean,lissack_clean,ibm_clean}.
Cammarata~\cite{cammarata_clean} утверждает, что видимая
полисемантичность нейронов в основном объясняется не суперпозицией, а
<<рекурсивным обобщением>> связанных понятий, и по предварительной
оценке около 90\% <<полюсов>> нейронов Qwen3-32B (положительная или отрицательная сторона активации) имеют одну основную функцию.
Gao~\cite{gao_notclean} возражает, что совпадения активирующих примеров с
шаблоном недостаточно, и предлагает строгий критерий: объяснение
верно, если подстановка смоделированных по нему активаций не
увеличивает KL-дивергенцию выходного распределения модели. Именно этот
критерий "--- KL при подстановке "--- мы используем как основную
метрику замены слоёв.

Эти наблюдения подсказывают иную постановку задачи. Если первые слои
LLM реализуют в основном локальные, <<поверхностные>> вычисления, то
их можно заменить \emph{явными детерминированными алгоритмами} вместо
единообразной аппроксимации внимания во всей сети. Под
детерминированным алгоритмом мы понимаем процедуру, которая
\begin{enumerate}
  \item не использует случайную выборку на этапе вывода;
  \item строится по весам обученной модели и/или статистике корпуса
    в замкнутом виде или явной процедурой, без градиентного
    дообучения;
  \item даёт одинаковый выход на одинаковом входе и допускает
    аналитическую оценку ошибки.
\end{enumerate}
Примеры таких алгоритмов: фиксированные разреженные шаблоны внимания,
позиционные ядра вместо голов внимания, таблицы поиска по токенам,
аналитические низкоранговые разложения матриц $W_{QK}$, $W_{OV}$ и
весов MLP~\cite{wong2025a3,yuan2023asvd,wang2025svdllm}, линейные
карты, подогнанные гребневой регрессией.

\paragraph{Постановка задачи.}
Пусть $f = f_L \circ \dots \circ f_1 \circ E$ "--- языковая модель из
$L$ блоков с матрицей эмбеддингов $E$, а $h_k(x)$ "--- состояние
остаточного потока после $k$-го блока на входе $x$. Требуется
построить детерминированный алгоритм $A_k$, такой что
\[
  \tilde f = f_L \circ \dots \circ f_{k+1} \circ A_k,
  \qquad
  \mathbb{E}_{x}\,
  \mathrm{KL}\bigl(p_f(\cdot \mid x)\,\|\,p_{\tilde f}(\cdot \mid x)\bigr)
  \le \varepsilon,
\]
и определить максимальную глубину $k^{*}(\varepsilon)$, для которой
это возможно при ограничениях на вычислительную сложность и память
алгоритма $A_k$.

\paragraph{Вклад работы.}
\begin{enumerate}
  \item Формализация задачи детерминированной аппроксимации первых
    $k$ слоёв LLM с критерием качества <<KL при подстановке>> и
    протокол, в котором исходная и изменённая модели сравниваются
    в одном процессе в fp32 (шум bf16 сам по себе даёт KL$\approx10^{-3}$).
  \item Адаптация к каузальному декодеру с GQA, QK-нормализацией и RoPE
    шести методов эффективного внимания и пяти методов сжатия MLP из
    литературы; разбор того, почему Linformer, Nyströmformer, FNet,
    SOFT и Reformer неприменимы без переобучения.
  \item Новые детерминированные методы: факторизация внимания
    <<позиционное ядро $\times$ значимость ключевого токена>> (O4),
    типизированная замена голов (O1), ридж-свёртка вместо
    внимания (O2), аддитивная $n$-граммная замена блоков (O3),
    линеаризация MLP с точными нейронами (M-O1), табличный (M-O2) и
    кусочно-линейный (M-O3) MLP. Все наши методы калибруются
    последовательно с компенсацией накопленного дрейфа.
  \item Послойный профиль Qwen3-0.6B и 1.7B, выделяющий режим <<до стока
    внимания>> и задающий глубину $k_{\mathrm{sink}}$, и анализ
    неприменимости случайных признаков к первым слоям.
  \item Сравнение по KL, perplexity и семи бенчмаркам на двух
    размерах модели, кривые <<качество -- бюджет вычислений>>,
    доверительные интервалы и дистилляционный контроль. Итог: внимание
    первых трёх слоёв заменяется окном со стоком за 1.5--1.8\,\% PPL без
    потерь на бенчмарках; MLP первых слоёв "--- лишь частично, лучший
    метод "--- наш предиктор активных нейронов M-O4 (4.7\,\% на 0.6B,
    1.7\,\% на 1.7B при четверти вычислений MLP).
  \item Ловушка калибровки в замкнутой форме: единичные токены с
    массивной активацией доминируют в задаче наименьших квадратов и
    обрушивают модель на слое, где возникает сток внимания. Мы даём
    формулу влияния строки, воспроизводим эффект дозой отравленных окон,
    устраняем его исключением строк-выбросов и показываем, какие
    опубликованные методы ему не подвержены и почему.
  \item Нейроны стока: массивную активацию создают один--пять нейронов
    MLP, без которых модель не работает; их сохранение "--- обязательное
    условие любого отбора нейронов в ранних слоях.
  \item Карта заменяемости по глубине и априорный критерий: дешёвые
    статистики слоя предсказывают KL его замены с ранговой корреляцией
    0.9--0.97, воспроизводятся на второй модели и позволяют планировать
    замены на непросмотренной модели в 10--15 раз точнее замены префикса.
\end{enumerate}

\paragraph{Структура статьи.}
Раздел~2 содержит обзор литературы. В разделе~3 описаны постановка,
адаптированные методы из литературы и предложенные методы. Раздел~4
описывает экспериментальную установку, раздел~5 "--- результаты, раздел~6
"--- обсуждение и ограничения.


% =====================================================================
% Общий список литературы (подключается через % Общий список литературы (подключается через % Общий список литературы (подключается через \input{references})
% Нумерация совпадает с literature.md
\begin{thebibliography}{99}
\selectlanguage{english}

\bibitem{elhage2021framework}
N.~Elhage et al. A Mathematical Framework for Transformer Circuits.
Transformer Circuits Thread, 2021.
\url{https://transformer-circuits.pub/2021/framework/index.html}

\bibitem{vaswani2017attention}
A.~Vaswani et al. Attention Is All You Need. NeurIPS, 2017.
arXiv:1706.03762.

\bibitem{geva2021ffn}
M.~Geva et al. Transformer Feed-Forward Layers Are Key-Value Memories.
EMNLP, 2021. arXiv:2012.14913.

\bibitem{cordonnier2020relationship}
J.-B.~Cordonnier et al. On the Relationship between Self-Attention and
Convolutional Layers. ICLR, 2020. arXiv:1911.03584.

\bibitem{dong2021attention}
Y.~Dong et al. Attention is Not All You Need: Pure Attention Loses Rank
Doubly Exponentially with Depth. ICML, 2021. arXiv:2103.03404.

\bibitem{cammarata_clean}
N.~Cammarata. Neural Networks Are Already Clean. Goodfire, 2026.
\url{https://already-clean.goodfire.com/}

\bibitem{gao_notclean}
L.~Gao. Neural networks are not already clean. LessWrong, 2026.
\url{https://www.lesswrong.com/posts/Dkh7pXKjoNWkkxQDz/neural-networks-are-not-already-clean}

\bibitem{olah2020zoom}
C.~Olah et al. Zoom In: An Introduction to Circuits. Distill, 2020.
\url{https://distill.pub/2020/circuits/zoom-in/}

\bibitem{lissack_clean}
M.~Lissack. Already Clean, Not Yet Meaningful: One ``how'' for the
question Cammarata has opened. PhilArchive.
\url{https://philarchive.org/rec/LISACN} % TODO: год

\bibitem{ibm_clean}
S.~Brodsky. Is AI easier to understand than scientists thought? IBM Think, 2026.
\url{https://www.ibm.com/think/news/is-ai-easier-to-understand-than-scientists-thought-goodfire-essay}

\bibitem{katharopoulos2020rnn}
A.~Katharopoulos et al. Transformers are RNNs: Fast Autoregressive
Transformers with Linear Attention. ICML, 2020. arXiv:2006.16236.

\bibitem{wang2020linformer}
S.~Wang et al. Linformer: Self-Attention with Linear Complexity, 2020.
arXiv:2006.04768.

\bibitem{xiong2021nystrom}
Y.~Xiong et al. Nyströmformer: A Nyström-Based Algorithm for
Approximating Self-Attention. AAAI, 2021. arXiv:2102.03902.

\bibitem{choromanski2021performer}
K.~Choromanski et al. Rethinking Attention with Performers. ICLR, 2021.
arXiv:2009.14794.

\bibitem{peng2021rfa}
H.~Peng et al. Random Feature Attention. ICLR, 2021. arXiv:2103.02143.

\bibitem{lu2021soft}
J.~Lu et al. SOFT: Softmax-free Transformer with Linear Complexity.
NeurIPS, 2021. arXiv:2110.11945.

\bibitem{child2019sparse}
R.~Child et al. Generating Long Sequences with Sparse Transformers, 2019.
arXiv:1904.10509.

\bibitem{kitaev2020reformer}
N.~Kitaev et al. Reformer: The Efficient Transformer. ICLR, 2020.
arXiv:2001.04451.

\bibitem{zaheer2020bigbird}
M.~Zaheer et al. Big Bird: Transformers for Longer Sequences. NeurIPS,
2020. arXiv:2007.14062.

\bibitem{leethorp2022fnet}
J.~Lee-Thorp et al. FNet: Mixing Tokens with Fourier Transforms.
NAACL, 2022. arXiv:2105.03824.

\bibitem{tay2021synthesizer}
Y.~Tay et al. Synthesizer: Rethinking Self-Attention in Transformer
Models. ICML, 2021. arXiv:2005.00743.

\bibitem{wong2025a3}
J.~T.~H.~Wong et al. A3: an Analytical Low-Rank Approximation Framework
for Attention, 2025. arXiv:2505.12942.

\bibitem{olsson2022induction}
C.~Olsson et al. In-context Learning and Induction Heads. Transformer
Circuits Thread, 2022. arXiv:2209.11895.

\bibitem{bricken2023monosemanticity}
T.~Bricken et al. Towards Monosemanticity: Decomposing Language Models
With Dictionary Learning. Transformer Circuits Thread, 2023.
\url{https://transformer-circuits.pub/2023/monosemantic-features/index.html}

\bibitem{rai2024review}
D.~Rai et al. A Practical Review of Mechanistic Interpretability for
Transformer-Based Language Models, 2024. arXiv:2407.02646.

\bibitem{yuan2023asvd}
Z.~Yuan et al. ASVD: Activation-aware Singular Value Decomposition for
Compressing Large Language Models, 2023. arXiv:2312.05821.

\bibitem{wang2025svdllm}
X.~Wang et al. SVD-LLM: Truncation-aware Singular Value Decomposition for
Large Language Model Compression. ICLR, 2025. arXiv:2403.07378.

\bibitem{liu2023dejavu}
Z.~Liu et al. Deja Vu: Contextual Sparsity for Efficient LLMs at Inference
Time. ICML, 2023. arXiv:2310.17157.

\bibitem{zhang2022moefication}
Z.~Zhang et al. MoEfication: Transformer Feed-forward Layers are Mixtures of
Experts. Findings of ACL, 2022. arXiv:2110.01786.

\bibitem{lee2024cats}
J.-Y.~Lee et al. CATS: Contextually-Aware Thresholding for Sparsity in Large
Language Models. COLM, 2024. arXiv:2404.08763.

\bibitem{liu2024teal}
J.~Liu et al. Training-Free Activation Sparsity in Large Language Models.
ICLR, 2025. arXiv:2408.14690.

\bibitem{shopkhoev2025replaceme}
D.~Shopkhoev et al. ReplaceMe: Network Simplification via Depth Pruning and
Transformer Block Linearization. NeurIPS, 2025. arXiv:2505.02819.

\bibitem{yun2026ghosted}
V.-D.~Yun et al. Ghosted Layers: Unconstrained Activation Alignment for
Recovering Layer-Pruned LLMs. ICML, 2026. arXiv:2605.15491.

\bibitem{yu2024superweights}
M.~Yu et al. The Super Weight in Large Language Models. arXiv:2411.07191,
2024.

\bibitem{sun2024massive}
M.~Sun et al. Massive Activations in Large Language Models. COLM, 2024.
arXiv:2402.17762.

\bibitem{merity2017pointer}
S.~Merity, C.~Xiong, J.~Bradbury, R.~Socher. Pointer Sentinel Mixture Models. ICLR, 2017.
arXiv:1609.07843. Набор WikiText: \url{https://huggingface.co/datasets/Salesforce/wikitext}

\bibitem{gao2020pile}
L.~Gao et al. The Pile: An 800GB Dataset of Diverse Text for Language Modeling. 2020.
arXiv:2101.00027. Подвыборка pile-10k: \url{https://huggingface.co/datasets/NeelNanda/pile-10k}

\bibitem{gao2023harness}
L.~Gao et al. A framework for few-shot language model evaluation (lm-evaluation-harness), v0.4. Zenodo, 2023.
\url{https://github.com/EleutherAI/lm-evaluation-harness}

\bibitem{paperno2016lambada}
D.~Paperno et al. The LAMBADA dataset: Word prediction requiring a broad discourse context. ACL, 2016.

\bibitem{zellers2019hellaswag}
R.~Zellers, A.~Holtzman, Y.~Bisk, A.~Farhadi, Y.~Choi. HellaSwag: Can a Machine Really Finish Your Sentence? ACL, 2019.

\bibitem{clark2018arc}
P.~Clark et al. Think you have Solved Question Answering? Try ARC, the AI2 Reasoning Challenge. 2018. arXiv:1803.05457.

\bibitem{bisk2020piqa}
Y.~Bisk, R.~Zellers, R.~Le~Bras, J.~Gao, Y.~Choi. PIQA: Reasoning about Physical Commonsense in Natural Language. AAAI, 2020.

\bibitem{sakaguchi2020winogrande}
K.~Sakaguchi, R.~Le~Bras, C.~Bhagavatula, Y.~Choi. WinoGrande: An Adversarial Winograd Schema Challenge at Scale. AAAI, 2020.

\bibitem{hendrycks2021mmlu}
D.~Hendrycks et al. Measuring Massive Multitask Language Understanding. ICLR, 2021. arXiv:2009.03300.

\end{thebibliography}
)
% Нумерация совпадает с literature.md
\begin{thebibliography}{99}
\selectlanguage{english}

\bibitem{elhage2021framework}
N.~Elhage et al. A Mathematical Framework for Transformer Circuits.
Transformer Circuits Thread, 2021.
\url{https://transformer-circuits.pub/2021/framework/index.html}

\bibitem{vaswani2017attention}
A.~Vaswani et al. Attention Is All You Need. NeurIPS, 2017.
arXiv:1706.03762.

\bibitem{geva2021ffn}
M.~Geva et al. Transformer Feed-Forward Layers Are Key-Value Memories.
EMNLP, 2021. arXiv:2012.14913.

\bibitem{cordonnier2020relationship}
J.-B.~Cordonnier et al. On the Relationship between Self-Attention and
Convolutional Layers. ICLR, 2020. arXiv:1911.03584.

\bibitem{dong2021attention}
Y.~Dong et al. Attention is Not All You Need: Pure Attention Loses Rank
Doubly Exponentially with Depth. ICML, 2021. arXiv:2103.03404.

\bibitem{cammarata_clean}
N.~Cammarata. Neural Networks Are Already Clean. Goodfire, 2026.
\url{https://already-clean.goodfire.com/}

\bibitem{gao_notclean}
L.~Gao. Neural networks are not already clean. LessWrong, 2026.
\url{https://www.lesswrong.com/posts/Dkh7pXKjoNWkkxQDz/neural-networks-are-not-already-clean}

\bibitem{olah2020zoom}
C.~Olah et al. Zoom In: An Introduction to Circuits. Distill, 2020.
\url{https://distill.pub/2020/circuits/zoom-in/}

\bibitem{lissack_clean}
M.~Lissack. Already Clean, Not Yet Meaningful: One ``how'' for the
question Cammarata has opened. PhilArchive.
\url{https://philarchive.org/rec/LISACN} % TODO: год

\bibitem{ibm_clean}
S.~Brodsky. Is AI easier to understand than scientists thought? IBM Think, 2026.
\url{https://www.ibm.com/think/news/is-ai-easier-to-understand-than-scientists-thought-goodfire-essay}

\bibitem{katharopoulos2020rnn}
A.~Katharopoulos et al. Transformers are RNNs: Fast Autoregressive
Transformers with Linear Attention. ICML, 2020. arXiv:2006.16236.

\bibitem{wang2020linformer}
S.~Wang et al. Linformer: Self-Attention with Linear Complexity, 2020.
arXiv:2006.04768.

\bibitem{xiong2021nystrom}
Y.~Xiong et al. Nyströmformer: A Nyström-Based Algorithm for
Approximating Self-Attention. AAAI, 2021. arXiv:2102.03902.

\bibitem{choromanski2021performer}
K.~Choromanski et al. Rethinking Attention with Performers. ICLR, 2021.
arXiv:2009.14794.

\bibitem{peng2021rfa}
H.~Peng et al. Random Feature Attention. ICLR, 2021. arXiv:2103.02143.

\bibitem{lu2021soft}
J.~Lu et al. SOFT: Softmax-free Transformer with Linear Complexity.
NeurIPS, 2021. arXiv:2110.11945.

\bibitem{child2019sparse}
R.~Child et al. Generating Long Sequences with Sparse Transformers, 2019.
arXiv:1904.10509.

\bibitem{kitaev2020reformer}
N.~Kitaev et al. Reformer: The Efficient Transformer. ICLR, 2020.
arXiv:2001.04451.

\bibitem{zaheer2020bigbird}
M.~Zaheer et al. Big Bird: Transformers for Longer Sequences. NeurIPS,
2020. arXiv:2007.14062.

\bibitem{leethorp2022fnet}
J.~Lee-Thorp et al. FNet: Mixing Tokens with Fourier Transforms.
NAACL, 2022. arXiv:2105.03824.

\bibitem{tay2021synthesizer}
Y.~Tay et al. Synthesizer: Rethinking Self-Attention in Transformer
Models. ICML, 2021. arXiv:2005.00743.

\bibitem{wong2025a3}
J.~T.~H.~Wong et al. A3: an Analytical Low-Rank Approximation Framework
for Attention, 2025. arXiv:2505.12942.

\bibitem{olsson2022induction}
C.~Olsson et al. In-context Learning and Induction Heads. Transformer
Circuits Thread, 2022. arXiv:2209.11895.

\bibitem{bricken2023monosemanticity}
T.~Bricken et al. Towards Monosemanticity: Decomposing Language Models
With Dictionary Learning. Transformer Circuits Thread, 2023.
\url{https://transformer-circuits.pub/2023/monosemantic-features/index.html}

\bibitem{rai2024review}
D.~Rai et al. A Practical Review of Mechanistic Interpretability for
Transformer-Based Language Models, 2024. arXiv:2407.02646.

\bibitem{yuan2023asvd}
Z.~Yuan et al. ASVD: Activation-aware Singular Value Decomposition for
Compressing Large Language Models, 2023. arXiv:2312.05821.

\bibitem{wang2025svdllm}
X.~Wang et al. SVD-LLM: Truncation-aware Singular Value Decomposition for
Large Language Model Compression. ICLR, 2025. arXiv:2403.07378.

\bibitem{liu2023dejavu}
Z.~Liu et al. Deja Vu: Contextual Sparsity for Efficient LLMs at Inference
Time. ICML, 2023. arXiv:2310.17157.

\bibitem{zhang2022moefication}
Z.~Zhang et al. MoEfication: Transformer Feed-forward Layers are Mixtures of
Experts. Findings of ACL, 2022. arXiv:2110.01786.

\bibitem{lee2024cats}
J.-Y.~Lee et al. CATS: Contextually-Aware Thresholding for Sparsity in Large
Language Models. COLM, 2024. arXiv:2404.08763.

\bibitem{liu2024teal}
J.~Liu et al. Training-Free Activation Sparsity in Large Language Models.
ICLR, 2025. arXiv:2408.14690.

\bibitem{shopkhoev2025replaceme}
D.~Shopkhoev et al. ReplaceMe: Network Simplification via Depth Pruning and
Transformer Block Linearization. NeurIPS, 2025. arXiv:2505.02819.

\bibitem{yun2026ghosted}
V.-D.~Yun et al. Ghosted Layers: Unconstrained Activation Alignment for
Recovering Layer-Pruned LLMs. ICML, 2026. arXiv:2605.15491.

\bibitem{yu2024superweights}
M.~Yu et al. The Super Weight in Large Language Models. arXiv:2411.07191,
2024.

\bibitem{sun2024massive}
M.~Sun et al. Massive Activations in Large Language Models. COLM, 2024.
arXiv:2402.17762.

\bibitem{merity2017pointer}
S.~Merity, C.~Xiong, J.~Bradbury, R.~Socher. Pointer Sentinel Mixture Models. ICLR, 2017.
arXiv:1609.07843. Набор WikiText: \url{https://huggingface.co/datasets/Salesforce/wikitext}

\bibitem{gao2020pile}
L.~Gao et al. The Pile: An 800GB Dataset of Diverse Text for Language Modeling. 2020.
arXiv:2101.00027. Подвыборка pile-10k: \url{https://huggingface.co/datasets/NeelNanda/pile-10k}

\bibitem{gao2023harness}
L.~Gao et al. A framework for few-shot language model evaluation (lm-evaluation-harness), v0.4. Zenodo, 2023.
\url{https://github.com/EleutherAI/lm-evaluation-harness}

\bibitem{paperno2016lambada}
D.~Paperno et al. The LAMBADA dataset: Word prediction requiring a broad discourse context. ACL, 2016.

\bibitem{zellers2019hellaswag}
R.~Zellers, A.~Holtzman, Y.~Bisk, A.~Farhadi, Y.~Choi. HellaSwag: Can a Machine Really Finish Your Sentence? ACL, 2019.

\bibitem{clark2018arc}
P.~Clark et al. Think you have Solved Question Answering? Try ARC, the AI2 Reasoning Challenge. 2018. arXiv:1803.05457.

\bibitem{bisk2020piqa}
Y.~Bisk, R.~Zellers, R.~Le~Bras, J.~Gao, Y.~Choi. PIQA: Reasoning about Physical Commonsense in Natural Language. AAAI, 2020.

\bibitem{sakaguchi2020winogrande}
K.~Sakaguchi, R.~Le~Bras, C.~Bhagavatula, Y.~Choi. WinoGrande: An Adversarial Winograd Schema Challenge at Scale. AAAI, 2020.

\bibitem{hendrycks2021mmlu}
D.~Hendrycks et al. Measuring Massive Multitask Language Understanding. ICLR, 2021. arXiv:2009.03300.

\end{thebibliography}
)
% Нумерация совпадает с literature.md
\begin{thebibliography}{99}
\selectlanguage{english}

\bibitem{elhage2021framework}
N.~Elhage et al. A Mathematical Framework for Transformer Circuits.
Transformer Circuits Thread, 2021.
\url{https://transformer-circuits.pub/2021/framework/index.html}

\bibitem{vaswani2017attention}
A.~Vaswani et al. Attention Is All You Need. NeurIPS, 2017.
arXiv:1706.03762.

\bibitem{geva2021ffn}
M.~Geva et al. Transformer Feed-Forward Layers Are Key-Value Memories.
EMNLP, 2021. arXiv:2012.14913.

\bibitem{cordonnier2020relationship}
J.-B.~Cordonnier et al. On the Relationship between Self-Attention and
Convolutional Layers. ICLR, 2020. arXiv:1911.03584.

\bibitem{dong2021attention}
Y.~Dong et al. Attention is Not All You Need: Pure Attention Loses Rank
Doubly Exponentially with Depth. ICML, 2021. arXiv:2103.03404.

\bibitem{cammarata_clean}
N.~Cammarata. Neural Networks Are Already Clean. Goodfire, 2026.
\url{https://already-clean.goodfire.com/}

\bibitem{gao_notclean}
L.~Gao. Neural networks are not already clean. LessWrong, 2026.
\url{https://www.lesswrong.com/posts/Dkh7pXKjoNWkkxQDz/neural-networks-are-not-already-clean}

\bibitem{olah2020zoom}
C.~Olah et al. Zoom In: An Introduction to Circuits. Distill, 2020.
\url{https://distill.pub/2020/circuits/zoom-in/}

\bibitem{lissack_clean}
M.~Lissack. Already Clean, Not Yet Meaningful: One ``how'' for the
question Cammarata has opened. PhilArchive.
\url{https://philarchive.org/rec/LISACN} % TODO: год

\bibitem{ibm_clean}
S.~Brodsky. Is AI easier to understand than scientists thought? IBM Think, 2026.
\url{https://www.ibm.com/think/news/is-ai-easier-to-understand-than-scientists-thought-goodfire-essay}

\bibitem{katharopoulos2020rnn}
A.~Katharopoulos et al. Transformers are RNNs: Fast Autoregressive
Transformers with Linear Attention. ICML, 2020. arXiv:2006.16236.

\bibitem{wang2020linformer}
S.~Wang et al. Linformer: Self-Attention with Linear Complexity, 2020.
arXiv:2006.04768.

\bibitem{xiong2021nystrom}
Y.~Xiong et al. Nyströmformer: A Nyström-Based Algorithm for
Approximating Self-Attention. AAAI, 2021. arXiv:2102.03902.

\bibitem{choromanski2021performer}
K.~Choromanski et al. Rethinking Attention with Performers. ICLR, 2021.
arXiv:2009.14794.

\bibitem{peng2021rfa}
H.~Peng et al. Random Feature Attention. ICLR, 2021. arXiv:2103.02143.

\bibitem{lu2021soft}
J.~Lu et al. SOFT: Softmax-free Transformer with Linear Complexity.
NeurIPS, 2021. arXiv:2110.11945.

\bibitem{child2019sparse}
R.~Child et al. Generating Long Sequences with Sparse Transformers, 2019.
arXiv:1904.10509.

\bibitem{kitaev2020reformer}
N.~Kitaev et al. Reformer: The Efficient Transformer. ICLR, 2020.
arXiv:2001.04451.

\bibitem{zaheer2020bigbird}
M.~Zaheer et al. Big Bird: Transformers for Longer Sequences. NeurIPS,
2020. arXiv:2007.14062.

\bibitem{leethorp2022fnet}
J.~Lee-Thorp et al. FNet: Mixing Tokens with Fourier Transforms.
NAACL, 2022. arXiv:2105.03824.

\bibitem{tay2021synthesizer}
Y.~Tay et al. Synthesizer: Rethinking Self-Attention in Transformer
Models. ICML, 2021. arXiv:2005.00743.

\bibitem{wong2025a3}
J.~T.~H.~Wong et al. A3: an Analytical Low-Rank Approximation Framework
for Attention, 2025. arXiv:2505.12942.

\bibitem{olsson2022induction}
C.~Olsson et al. In-context Learning and Induction Heads. Transformer
Circuits Thread, 2022. arXiv:2209.11895.

\bibitem{bricken2023monosemanticity}
T.~Bricken et al. Towards Monosemanticity: Decomposing Language Models
With Dictionary Learning. Transformer Circuits Thread, 2023.
\url{https://transformer-circuits.pub/2023/monosemantic-features/index.html}

\bibitem{rai2024review}
D.~Rai et al. A Practical Review of Mechanistic Interpretability for
Transformer-Based Language Models, 2024. arXiv:2407.02646.

\bibitem{yuan2023asvd}
Z.~Yuan et al. ASVD: Activation-aware Singular Value Decomposition for
Compressing Large Language Models, 2023. arXiv:2312.05821.

\bibitem{wang2025svdllm}
X.~Wang et al. SVD-LLM: Truncation-aware Singular Value Decomposition for
Large Language Model Compression. ICLR, 2025. arXiv:2403.07378.

\bibitem{liu2023dejavu}
Z.~Liu et al. Deja Vu: Contextual Sparsity for Efficient LLMs at Inference
Time. ICML, 2023. arXiv:2310.17157.

\bibitem{zhang2022moefication}
Z.~Zhang et al. MoEfication: Transformer Feed-forward Layers are Mixtures of
Experts. Findings of ACL, 2022. arXiv:2110.01786.

\bibitem{lee2024cats}
J.-Y.~Lee et al. CATS: Contextually-Aware Thresholding for Sparsity in Large
Language Models. COLM, 2024. arXiv:2404.08763.

\bibitem{liu2024teal}
J.~Liu et al. Training-Free Activation Sparsity in Large Language Models.
ICLR, 2025. arXiv:2408.14690.

\bibitem{shopkhoev2025replaceme}
D.~Shopkhoev et al. ReplaceMe: Network Simplification via Depth Pruning and
Transformer Block Linearization. NeurIPS, 2025. arXiv:2505.02819.

\bibitem{yun2026ghosted}
V.-D.~Yun et al. Ghosted Layers: Unconstrained Activation Alignment for
Recovering Layer-Pruned LLMs. ICML, 2026. arXiv:2605.15491.

\bibitem{yu2024superweights}
M.~Yu et al. The Super Weight in Large Language Models. arXiv:2411.07191,
2024.

\bibitem{sun2024massive}
M.~Sun et al. Massive Activations in Large Language Models. COLM, 2024.
arXiv:2402.17762.

\bibitem{merity2017pointer}
S.~Merity, C.~Xiong, J.~Bradbury, R.~Socher. Pointer Sentinel Mixture Models. ICLR, 2017.
arXiv:1609.07843. Набор WikiText: \url{https://huggingface.co/datasets/Salesforce/wikitext}

\bibitem{gao2020pile}
L.~Gao et al. The Pile: An 800GB Dataset of Diverse Text for Language Modeling. 2020.
arXiv:2101.00027. Подвыборка pile-10k: \url{https://huggingface.co/datasets/NeelNanda/pile-10k}

\bibitem{gao2023harness}
L.~Gao et al. A framework for few-shot language model evaluation (lm-evaluation-harness), v0.4. Zenodo, 2023.
\url{https://github.com/EleutherAI/lm-evaluation-harness}

\bibitem{paperno2016lambada}
D.~Paperno et al. The LAMBADA dataset: Word prediction requiring a broad discourse context. ACL, 2016.

\bibitem{zellers2019hellaswag}
R.~Zellers, A.~Holtzman, Y.~Bisk, A.~Farhadi, Y.~Choi. HellaSwag: Can a Machine Really Finish Your Sentence? ACL, 2019.

\bibitem{clark2018arc}
P.~Clark et al. Think you have Solved Question Answering? Try ARC, the AI2 Reasoning Challenge. 2018. arXiv:1803.05457.

\bibitem{bisk2020piqa}
Y.~Bisk, R.~Zellers, R.~Le~Bras, J.~Gao, Y.~Choi. PIQA: Reasoning about Physical Commonsense in Natural Language. AAAI, 2020.

\bibitem{sakaguchi2020winogrande}
K.~Sakaguchi, R.~Le~Bras, C.~Bhagavatula, Y.~Choi. WinoGrande: An Adversarial Winograd Schema Challenge at Scale. AAAI, 2020.

\bibitem{hendrycks2021mmlu}
D.~Hendrycks et al. Measuring Massive Multitask Language Understanding. ICLR, 2021. arXiv:2009.03300.

\end{thebibliography}


\end{document}
